%%%%%%%%%%%%%%%%%%%%%%%%%%%%%%%%%%%%%%%%%%%%%%%%
\usepackage[a4paper]{geometry}		
%
%% for example, change the margins to 2 inches all round
%%\geometry{top=1.25cm, bottom=-.6cm, left=1.5cm, right=1.5cm}
\geometry{top=2.540cm, bottom=2.540cm, left=2.540cm, right=2.540cm} 



% \usepackage[section]{placeins}
\setlength{\textfloatsep}{8pt plus 2pt minus 2pt}
%%% 8-12pt
%%%%%%%%%%%%%%%%%%%55
\usepackage{amsmath,amsfonts,amssymb}
% \usepackage{mathrsfs} % \mathscr
\usepackage[scr=rsfs]{mathalpha} % \mathscr
\usepackage{mathtools}
\usepackage{dsfont} % \mathds{1}
\usepackage{bm}

%\usepackage[numbers,sort]{natbib} # author year
\usepackage{xurl,xcolor,enumerate}
% \usepackage[nocompress,sort]{cite} %for numbers

%%%% table setting
\usepackage{tabularx,multirow,array,booktabs}
\usepackage{colortbl}%%% color in tables
%%%%%%%%%%%%%%%%%%%%%%%%%%%%%%%%%%%%%%%%%%%%%%%%%%%%
%%%%%% figure setting
\usepackage{graphicx}%% 
\graphicspath{{figures/}}%
% \graphicspath{{figures/}{figuresOld/}}%
\usepackage{float}%% 
\usepackage{subcaption} %% subfigure env
% \captionsetup{subrefformat=parens} 
\renewcommand\thesubfigure{(\alph{subfigure})}
\captionsetup[subfigure]{labelformat=simple, labelsep=space, 
% labelfont=bf
}
% \usepackage{subcaption} %% subfigure env
% \captionsetup{subrefformat=parens} 
\captionsetup[subfigure]{aboveskip=3pt}
\usepackage{caption}
\captionsetup[figure]{aboveskip=7pt}
\captionsetup[figure]{belowskip=2pt}
\captionsetup[table]{aboveskip=5pt}
%%%%%%%%%%%%%%%%%%%%%%%%%%%%%%%%%%
%%%%%%%%%%%%%%%%%%%%%%%%%%%%%%%%%%%%%
%\makeatletter\@addtoreset{equation}{section}\makeatother
%%% \thequation resets when \thesection updates
%\renewcommand{\theequation}{\arabic{section}.\arabic{equation}}
% %%%%%%%%%%%%%%%%%%%%%%%%%%%%%%%%%
% %%%%%%% set page properties
% \setlength{\oddsidemargin}{0.5cm}
% \setlength{\evensidemargin}{0.5cm}
% \setlength{\textwidth}{15.8cm}
% \setlength{\topmargin}{0.4cm}
% \setlength{\headheight}{0.0cm}% 
% \setlength{\headsep}{0.0cm}
% \setlength{\textheight}{23.5cm}
% Specify the dimensions of each page
%%%%%%%%%%%%%%%%%%%%%%%%%%%%%%%
%%%%%%%%%%%%%%%%%%%%%%%%%%%%%%%%%%
%% Theorem environment below
\usepackage{amsthm}
\theoremstyle{plain}
\newtheorem{theorem}{Theorem}[section]%  
\newtheorem{corollary}[theorem]{Corollary} %
\newtheorem*{main}{Main Theorem}% 
\newtheorem{lemma}[theorem]{Lemma}
\newtheorem*{lemma*}{Lemma}
\newtheorem{problem}[theorem]{Problem}
\newtheorem{claim}[theorem]{Claim}
\newtheorem{proposition}[theorem]{Proposition}
\theoremstyle{definition}
\newtheorem{definition}[theorem]{Definition}
\theoremstyle{remark}
\newtheorem*{notation}{Notation}
\newtheorem{remark}[theorem]{Remark}
\newtheorem*{remark*}{Remark}
%%%%%%%%%%%%%%%%%%%%%%%%%%%%%%%%%%%%%%%%%%%%%%%%%%%%%%%%%%%%%%%%%%%%%%%%%%%%
%%%% New commands below
%%%% math fonts
\usepackage{etoolbox,xparse,dsfont,bm,amssymb}
\newcommand{\myMFabc}[4]{\expandafter#1\csname#3#4\endcsname{{#2{#4}}}}
%\forcsvlist{\myMF{\newcommand}{\bm}{bm}}{x,y}
\newcommand{\myMFcmd}[4]{\expandafter#1\csname#3#4\endcsname{{#2{\csname#4\endcsname}}}}
%\forcsvlist{\myMFcmd{\newcommand}{\bm}{bm}}{theta,phi}


\newcommand{\MFabc}[3][\newcommand]{
    \def\doOld##1##2{\forcsvlist{\myMFabc{#1}{##1}{##2}}{#3}}
    \providecommand{\do}{do}
    \RenewDocumentCommand \do { >{\SplitList{,}} m } { \doOld##1 }
    \docsvlist{#2}
}
% \MFabc{ {\bm,bm}, {\mathcal,cal} }{A,B,C}
\newcommand{\MFcmd}[3][\newcommand]{
    \def\doOld##1##2{\forcsvlist{\myMFcmd{#1}{##1}{##2}}{#3}}
    \providecommand{\do}{do}
    \RenewDocumentCommand \do { >{\SplitList{,}} m } { \doOld##1 }
    \docsvlist{#2}
}

%%% 0 and 1
\newcommand{\bmzero}{{\bm{0}}}
\newcommand{\bmone}{{\bm{1}}}
\def\bbone{{\ensuremath{\mathds{1}}}}
\let\one\bbone
%  mathfrak, e.g., \frakT = \mathfrak{T}
\MFabc{ {\mathfrak,frak}, }{T,u,U,o,O,E,m}
%  mathbb, e.g., \R = \mathbb{R}
\MFabc{ {\mathbb, }, }{A,N,Z,R,C,Q}

%%% all capital letters
% \mathcal, \mathscr, and \mathbb 
% e.g., \calA=\mathcal{A}
% \bm\mathcal, e.g., \calbmA = \bm{\mathcal{A}}
\newcommand{\hatbm}[1]{\widehat{\bm{#1}}}
\newcommand{\tildebm}[1]{\widetilde{\bm{#1}}}
\newcommand{\bmcal}[1]{\bm{\mathcal{#1}}}
\newcommand{\caltilde}[1]{\mathcal{\widetilde{#1}}}
\newcommand{\calhat}[1]{\mathcal{\widehat{#1}}} 
\newcommand{\scrtilde}[1]{\widetilde{\mathscr{#1}\mspace{1mu}\mspace{-1mu}}}% \mspace{1mu}\mspace{-1mu} to avoid spacing conflict \mathscr and \widetilde
\newcommand{\scrhat}[1]{\mathscr{\widehat{#1}\mspace{1mu}\mspace{-1mu}}}
\newcommand{\bmcalhat}[1]{\bm{\mathcal{\widehat{#1}}}}
\newcommand{\bmcaltilde}[1]{\bm{\mathcal{\widetilde{#1}}}}
%%%%%%
\MFabc{ {\mathcal,cal}, {\mathscr,scr}, {\mathbb,bb}, {\bmcal,bmcal}, {\bmcal,calbm}, {\caltilde,caltilde}, {\caltilde,tildecal}, {\calhat,calhat}, {\calhat,hatcal}, {\scrtilde,scrtilde}, {\scrtilde,tildescr}, {\scrhat,scrhat}, {\scrhat,hatscr}, {\bmcalhat,bmcalhat}, {\bmcalhat,bmhatcal}, 
{\bmcalhat,calbmhat}, {\bmcalhat,calhatbm}, {\bmcalhat,hatbmcal}, {\bmcalhat,hatcalbm}, {\bmcaltilde,bmcaltilde}, {\bmcaltilde,bmtildecal}, {\bmcaltilde,calbmtilde}, {\bmcaltilde,caltildebm}, {\bmcaltilde,tildebmcal}, {\bmcaltilde,tildecalbm}}{A,B,C,D,E,F,G,H,I,J,K,L,M,N,O,P,Q,R,S,T,U,V,W,X,Y,Z}

\renewcommand{\scrtildeA}{\widetilde{\mathscr{A}\mspace{2mu}}\mspace{-6.1mu}}% \mspace{2mu} and\mspace{-6.1mu} to avoid spacing conflict \mathscr and \widetilde
\renewcommand{\scrhatA}{\mathscr{\widehat{A}\mspace{2mu}}\mspace{-6.1mu}}
\let\tildescrA\scrtildeA
\let\hatscrA\scrhatA

%%% all letters
% bm, \widehat， \widetilde, and
% \mathsf shortcuts -- for random variables， 
% e.g., \bmA=\bm{A}, \bmalpha=\bm{\alpha}
%%%%%%%%%
\MFabc{{\bm,bm}, {\mathsf,sf}, {\widehat,hat}, {\widetilde,tilde}, {\hatbm,hatbm}, {\hatbm,bmhat}, {\tildebm,tildebm}, {\tildebm,bmtilde}}{a,b,c,d,e,f,g,h,i,j,k,l,m,n,o,p,q,r,s,t,u,v,w,x,y,z,A,B,C,D,E,F,G,H,I,J,K,L,M,N,O,P,Q,R,S,T,U,V,W,X,Y,Z}

%%% all Greek alphabet, e.g., \bmalpha = \bm{\alpha}
\let\eps\varepsilon
\MFcmd{{\bm,bm}, {\widehat,hat}, {\widetilde,tilde}, {\tildebm, tildebm}, {\tildebm, bmtilde}, {\hatbm, hatbm}, {\hatbm, bmhat}}{alpha,beta,gamma,delta,epsilon,zeta,eta,theta,iota,kappa,lambda,mu,nu,xi,omicron,pi,rho,sigma,tau,upsilon,phi,chi,psi,omega,Alpha,Beta,Gamma,Delta,Epsilon,Zeta,Eta,Theta,Iota,Kappa,Lambda,Mu,Nu,Xi,Omicron,Pi,Rho,Sigma,Tau,Upsilon,Phi,Chi,Psi,Omega,varrho,varphi,vartheta,varepsilon,varsigma,ell,eps}



%%%%%%%%%% activation functions
%\newcommand{\actdef}[1]{\lowercase{\expandafter\def\csname#1\endcsname}{\texttt{#1}}}
\newcommand{\actdef}[1]{\expandafter\def\csname#1\endcsname{{\ensuremath{\mathtt{#1}}}}}
\forcsvlist{\actdef}{ReLU, LReLU, LeakyReLU, ELU, GELU, SiLU, Softplus, dGELU, dSiLU, dSoftplus, Tanh, Sigmoid, Arctan, Softsign, SRS, dSRS, Swish, dSwish, Mish, dMish, SELU, CELU, dSELU, Sin,SinLU, SinTU, PSinTU, sine, cosine, Sine, Cosine, EUAF}

% \newcommand{\actdefs}[1]{\expandafter\def\csname#1s\endcsname{{\ensuremath{\mathtt{#1}}s}}}
% \forcsvlist{\actdefs}{ELU, GELU, SiLU, dGELU, dSiLU, SRS, dSRS, dSwish}


%%%%%%%%%%%%%%%%%%%%%%%%%%%%%%%%%%%%%%%%%%

%\newcommand{\nn}{{\hspace{0.6pt}\mathcal{N}\hspace{-1.80pt}\mathcal{N}\hspace{-1.3pt}}}
% \newcommand{\nn}[6][]{\ensuremath{
% 		{\hspace{0.6pt}\mathcal{N}\hspace{-1.98pt}\mathcal{N}\hspace{-0.725pt}}_{\hspace{-0.75pt}#2\hspace{-0.02pt}}%  NN
% 		#1\{#3,\hspace{1.987pt} #4;\hspace{3.0297pt} \R^{#5}\hspace{-1.0298pt}\to\hspace{-0.98pt}\R^{#6}#1\}
% }}
% \newcommand{\mmnn}{\ensuremath{	{\hspace{0.6pt}\mathcal{M}\hspace{-1.98pt}\mathcal{M}\hspace{-1.98pt}\mathcal{N}\hspace{-1.98pt}\mathcal{N}\hspace{-0.725pt}}
% }}

\newcommand{\mn}[7][]{\ensuremath{	{\hspace{0.6pt}\mathcal{M}\hspace{-0.0pt}\mathcal{N}\hspace{-0.8pt}}_{#2}%  NN
		#1\{#3,\hspace{2.2pt} #4,\hspace{2.2pt} #5;\hspace{3.0297pt} {#6}\hspace{-1.0298pt}\to\hspace{-0.98pt}{#7}#1\}
}}


\let\mmnn\mn

\newcommand{\nn}[6][]{\ensuremath{	{\hspace{0.6pt}\mathcal{N}\hspace{-1.28pt}\mathcal{N}\hspace{-0.725pt}}_{#2}%  NN
		#1\{#3,\hspace{1.987pt} #4;\hspace{3.0297pt} {#5}\hspace{-1.0298pt}\to\hspace{-0.98pt}{#6}#1\}
}}

\renewcommand{\nn}[6][]{\ensuremath{	{\hspace{0.6pt}\mathcal{F}\hspace{-0.828pt}\mathcal{N}\hspace{-0.725pt}}_{#2}%  NN
		#1\{#3,\hspace{1.987pt} #4;\hspace{3.0297pt} {#5}\hspace{-1.0298pt}\to\hspace{-0.98pt}{#6}#1\}
}}
% \let\nn\mn
\let\fcnn\nn

% \NewDocumentCommand{\nn}{o o m m m m}{\ensuremath{		{\hspace{0.6pt}\mathcal{N}\hspace{-1.98pt}\mathcal{N}\hspace{-0.725pt}}_{\hspace{-0.75pt}#2\hspace{-0.02pt}}%  NN
%     		#1\{#3,\hspace{1.987pt} #4;\hspace{3.0297pt} \R^{#5}\hspace{-1.0298pt}\to\hspace{-0.98pt}\R^{#6}#1\}
% }}

% \NewDocumentCommand{\nnOneD}{o o m m m m}{\ensuremath{
% {\hspace{0.6pt}\mathcal{N}\hspace{-1.98pt}\mathcal{N}\hspace{-0.725pt}}_{#2}%  NN
% 		#1\{#3,\hspace{1.987pt} #4;\hspace{3.0297pt} {#5}\hspace{-1.0298pt}\to\hspace{-0.98pt}{#6}#1\}
% }}


%\newcommand{\tbinom}[2]{{\textstyle\binom{#1}{#2}}}
% \newcommand{\ts}{{T}}
\let\ts\intercal
\def\ts{\tn{$\textsf{T}$}}
\def\ts{\textsf{\textup{T}}}

% \newcommand{\cpwl}{\tn{CPwL}}
\def\cpwl{\textsf{\textup{CPwL}}}
\let\cpl\cpwl
\let\CPwL\cpwl
\newcommand{\din}{{d_\tn{in}}}
\newcommand{\dout}{{d_\tn{out}}}
\newcommand{\dprime}{{\prime\prime}}
\newcommand{\tprime}{{\prime\prime\prime}}

%%%%%%%%%%%%%%%%%%%%%%%%%%%%%%%%%%%%%%%%%%%%%%%%%%%%%%%%%
\DeclareMathOperator*{\argmax}{arg\,max}
\DeclareMathOperator*{\argmin}{arg\,min}


%%% see web http://phaseportrait.blogspot.com/2007/08/lineno-and-amsmath-compatibility.html
%%%%line number below
\usepackage[left,mathlines]{lineno}
\usepackage{refcount}
%\renewcommand\linenumberfont{\rmfamily\upshape\scriptsize}

%%%%%%%%%%%%%%%%%%%%%%%%%%%%%%%%%%%%%%%%%%
%%%%%% %%%%% line numbers (cross-ref) properties setting
\renewcommand\thelinenumber{{\rmfamily\upshape\normalsize\color{blue}\arabic{linenumber}}}
%%%%%%%%%%%%%%%%%%%%%%%%%%%%%%%%%%%%%%%%%
%%%%% line numbers properties setting
\definecolor{mylinenumbercolor}{HTML}{BEBEBE}
\renewcommand\LineNumber{{\rmfamily\upshape\footnotesize\color{mylinenumbercolor}\arabic{linenumber}}}
\linenumbers
%%%%%%%%%%%%%%%%%%%%%%%%%%%%%%%%%%%%%%%%%%
\newcommand{\mylinelabel}[1]{\phantomsection\label{#1zsj}\linelabel{#1}}
%\newcommand{\mylineref}[1]{\hyperref[#1zsj]{\ref{#1}}}
%%%%%%%%% modify line numbers
\makeatletter
\newcommand*\patchAmsMathEnvironmentForLineno[1]{%
	\expandafter\let\csname old#1\expandafter\endcsname\csname #1\endcsname
	\expandafter\let\csname oldend#1\expandafter\endcsname\csname end#1\endcsname
	\renewenvironment{#1}%
	{\linenomath\csname old#1\endcsname}%
	{\csname oldend#1\endcsname\endlinenomath}}% 
\newcommand*\patchBothAmsMathEnvironmentsForLineno[1]{%
	\patchAmsMathEnvironmentForLineno{#1}%
	\patchAmsMathEnvironmentForLineno{#1*}}%
\patchBothAmsMathEnvironmentsForLineno{equation}%
\patchBothAmsMathEnvironmentsForLineno{align}%
\patchBothAmsMathEnvironmentsForLineno{flalign}%
\patchBothAmsMathEnvironmentsForLineno{alignat}%
\patchBothAmsMathEnvironmentsForLineno{gather}%
\patchBothAmsMathEnvironmentsForLineno{multline}%
\makeatother
%%%%%%%%%%%%%%%%%%%%%%%%%%%%%%%%%%%%%%%%%%%%%%%%%%%%%%%%

\newcommand{\mailto}[2][]{\href{mailto:#2?cc=#1}{\color{black}#2}}



% \usepackage{doi}\def\doitext{doi: }
%%%%%%%%%%%%%%%%%%%%%%%%%%
\let\tilde\widetilde
\let\hat\widehat
\let\epsilon\varepsilon
\let\eps\varepsilon
\let\subset\subseteq
\let\tn\textnormal
%%%% 
% \let\mycdots\cdots
% \def\cdots{\mathop{\mycdots}}
\newcommand{\customcdots}{%
	\mathinner{\mkern-0.1mu\cdotp\mkern-0.3mu\cdotp\mkern-0.3mu\cdotp\mkern-0.1mu}%
}
\let\cdots\customcdots
\let\myforall\forall
\def\forall{{\myforall\, }}
\let\myexists\exists
\def\exists{{\myexists\, }}
\let\emptyset\varnothing
%%%%
\long\def\black#1{{\color{black}#1}}
\long\def\red#1{{\color{red}#1}}
\long\def\green#1{{\color{green}#1}}
\long\def\blue#1{{\color{blue}#1}}
\long\def\cyan#1{{\color{cyan}#1}}

% define colors used in table
\definecolor{mygray}{RGB}{230,230,230}
\definecolor{myorange}{HTML}{ff7f0e}

%%%%%%%%%%%%%%%%%%%%%%%%%%%%%%%%%

\let\cite\citep
%%%\clearpage
% \bibliographystyle{siam}
% \bibliographystyle{plainnat}
% \bibliographystyle{plain}	
% \bibliographystyle{plainnat}
\usepackage{doi}
\def\doitext{DOI: }
% \def\url#1{\href{#1}{\nolinkurl{#1}}}
% \urlstyle{sf}



%%%%%%%%%%%%%%%%%%%%%%%%%%%%%%%%%%%%%%%%%%%%%%%%%%%%%%%%
%%%% define title page, following JMLR template
\makeatletter
\long\def\@makefntext#1{\@setpar{\@@par\@tempdima \hsize 
		\advance\@tempdima-15pt\parshape \@ne 15pt \@tempdima}\par
	\parindent 2em\noindent \hbox to \z@{\hss{\textsuperscript{\@thefnmark}} \hfil}#1}
%%%%%%%%%%%%%%%%%
\newlength\aftertitskip     \newlength\beforetitskip
\newlength\interauthorskip  \newlength\aftermaketitskip
%% Changeable parameters.
\setlength\aftertitskip{0.1in plus 0.2in minus 0.2in}
\setlength\beforetitskip{0.05in plus 0.08in minus 0.08in}
\setlength\interauthorskip{0.08in plus 0.1in minus 0.1in}
\setlength\aftermaketitskip{0.3in plus 0.1in minus 0.1in}
%%%%%%%%%%%%%%
%% overall definition of maketitle, @maketitle does the real work
\def\maketitle{\par
	\begingroup
	\def\thefootnote{\color{black}\fnsymbol{footnote}}
	\def\@makefnmark{\hbox to 0pt{$^{\@thefnmark}$\hss}}
	\@maketitle \@thanks
	\endgroup
	\setcounter{footnote}{0}
	\let\maketitle\relax \let\@maketitle\relax
	\gdef\@thanks{}\gdef\@author{}\gdef\@title{}\let\thanks\relax}

\def\@startauthor{\noindent \normalsize\bf}
\def\@endauthor{}
\def\@starteditor{\noindent \small {\bf Editor:~}}
\def\@endeditor{\normalsize}
\def\@maketitle{\vbox{\hsize\textwidth
		\linewidth\hsize \vskip \beforetitskip
		{\begin{center} \Large\bf \@title \par \end{center}} \vskip \aftertitskip
		{\def\and{\unskip\enspace{\rm and}\enspace}%
			\def\addr{\small\it}%
			\def\email{\hfill\small\sc}%
            % \def\email{\hfill\small}%
			\def\name{\normalsize\bf}%
			\def\AND{\@endauthor\rm\hss \vskip \interauthorskip \@startauthor}
			\@startauthor \@author \@endauthor}
		\vskip \aftermaketitskip
}}
\makeatother
% %%%%%%%%%%%%% title page setting
\addtolength\aftertitskip{0.25in}
\addtolength\aftermaketitskip{0.12in}
\addtolength\interauthorskip{4pt}
% %%%%%%%%%%%%%%%%%%%%




% === Float Placement Settings ===
\renewcommand{\topfraction}{0.95}       
% Allows floats (figures/tables) to occupy up to 95% of the top portion of a page.
% Default is usually 0.7. Increasing this makes it easier for large figures to appear at the top.

\renewcommand{\bottomfraction}{0.9}

\renewcommand{\textfraction}{0.05}      
% Requires that at least 5% of a page must be non-float text.
% Default is 0.2. Lowering this makes it easier to place large floats on the same page with minimal text.

\renewcommand{\floatpagefraction}{0.95}  
% A float-only page (i.e., a page with only figures or tables) must be at least 90% full.
% Default is 0.5. Increasing this helps prevent small floats from occupying a whole page.

\setcounter{topnumber}{5}     
% Maximum number of floats allowed at the top of a page.
% Default is usually 2.

\setcounter{bottomnumber}{5}  
% Maximum number of floats allowed at the bottom of a page.
% Default is usually 1.

\setcounter{totalnumber}{10}  
% Maximum total number of floats allowed on a single page.
% Default is usually 3.

% These settings help encourage LaTeX to place figures and tables near their references,
% and reduce the chance of floats getting pushed to separate pages.